\documentclass[11pt]{ctexart}

\usepackage{amsmath,amssymb}
\usepackage{booktabs,tabularx,array}
\usepackage[margin=2.2cm]{geometry}
\usepackage[authoryear,round]{natbib}
\usepackage{xcolor}
\usepackage{hyperref}

\newcommand{\tbv}{\textcolor{red}{[待核]}}
\newcolumntype{Y}{>{\raggedright\arraybackslash}X}
\setcounter{section}{2}

\begin{document}

\section{随机加速的研究格局}

比较随机加速的结果，必须同时固定四件事：噪声模型、取值点、收敛意义与总样本预算。
几乎必然（a.s.）意义下还要区分“速率”与“收敛”：
\[
  f(x_k)-f^\star=O(r_k)\quad\text{a.s.}
\]
表示对几乎每条轨道存在（随机）常数 $C(\omega)$，使 $f(x_k)-f^\star\le C(\omega)\,r_k$；而
\[
  f(x_k)\to f^\star\quad\text{a.s.}
\]
只说明误差趋于零，不给出任何速率。按噪声体制，现有文献可分为五条子线（表~\ref{tab:lines}）。

\begin{table}[htbp]
\centering
\small
\caption{随机加速的五条子线（凸、$L$-光滑）}
\label{tab:lines}
\begin{tabularx}{\textwidth}{@{}p{2.1cm}p{3.0cm}Yp{3.0cm}@{}}
\toprule
子线 & 噪声体制 & 已知结果 & 代表工作 \\
\midrule
期望复杂度
& 持续加性噪声，方差有界
& 已知总迭代数 $T$ 时，加速方法达到
  $O\!\big(LD^2/T^2+\sigma D/\sqrt T\big)$，与下界
  $\Omega\!\big(LD^2/T^2+\sigma D/\sqrt T\big)$ 匹配；
  加速只改善偏差项。不依赖 $T$ 的版本差一个对数因子。
& \citet{lan2012,cutkosky2019anytime}；
  下界 \citet{nemirovski1983,agarwal2012} \\
\addlinespace
动量的 a.s.\ 理论
& 持续加性噪声
& SHB，动量由步长决定（约为 $1-(3-4\varepsilon)/k$）：
  末次点 $o(k^{-1/2+\varepsilon})$ a.s.，迭代点 a.s.\ 收敛；
  SHB、SNAG，常数动量：$O(k^{-1/3+\varepsilon})$ a.s.，迭代点 a.s.\ 收敛。
& \citet{sebbouh2021,liu2022} \\
\addlinespace
相对噪声
& $\mathbb E[\|e_k\|^2\mid\mathcal F_k]\le\nu^2\|\nabla f(y_k)\|^2$
& 期望意义下可恢复 $O(1/k^2)$：原始 NAG 需要 $\nu<1$，
  修正方法可适用于任意 $\nu$；标准参数下一般不加速，甚至可能发散；
  a.s.\ 速率与迭代点收敛\tbv。
& \citet{vaswani2019fast,gupta2024nesterov,yu2026shang}；
  负结果 \citet{kidambi2018insufficiency,liu2020accelerating} \\
\addlinespace
噪声递减
& $\sigma_k\to0$，满足 $(K_2)$ 或 $(K_2^+)$
& 一般 $\alpha_k$ 的 SNAG/SRAG：a.s.\ $O\big(1/(s_kt_k^2)\big)$ 及其 $o$ 形式、
  梯度可和、迭代点 a.s.\ 收敛；临界参数下 FISTA 的迭代点 a.s.\ 弱收敛（条件\tbv）。
& \citet{attouch2026ravine,salzo2025iterates} \\
\addlinespace
连续时间随机动力学
& 随机惯性 SDE，含消失阻尼或 Hessian 阻尼
& 期望与 a.s.\ 意义下的收敛性质；用时间尺度变换与平均，
  把一阶随机动力学的结论转移到二阶系统。
& \citet{maulensoto2026stochastic,maulensoto2025sde} \\
\bottomrule
\end{tabularx}
\end{table}

\paragraph{持续噪声：主导阶被下界锁定。}
前两条子线说明，持续噪声下任何方法的噪声项都不快于 $\sigma D/\sqrt T$，加速只能改善偏差项。
a.s.\ 意义下，动量方法末次点的最好速率任意接近 $k^{-1/2}$（从慢的一侧逼近）。
动量由步长决定的 SHB 达到这一水平，常数动量只到 $k^{-1/3}$；
差别在于前者的末次点是全历史平均，后者只是短记忆平均。

\paragraph{相对噪声：加速真实存在，但原始 SNAG 并不总是有效。}
条件写在 $y_k$ 上，因为 SNAG 在外推点 $y_k$ 处计算随机梯度。
\citet{vaswani2019fast} 使用二阶矩形式的强增长条件
$\mathbb E\|g_k\|^2\le\rho\|\nabla f\|^2$；在无偏假设下，它等价于
$\mathbb E\|e_k\|^2\le(\rho-1)\|\nabla f\|^2$。
负结果同样重要：标准参数的 Nesterov SGD 在插值设定下一般不能保证加速，甚至可能发散
\citep{kidambi2018insufficiency,liu2020accelerating}。
因此问题不是“原始 SNAG 在相对噪声下是否有效”，而是“在什么噪声比例和参数下有效”。
SHANG++ \citep{yu2026shang} 由 Hessian 驱动的 Nesterov 流离散得到，
而 Ravine 方法的高分辨率 ODE 同样含 Hessian 驱动阻尼 \citep{attouch2026ravine}。
\emph{推断：}相对噪声下的稳健加速可能依赖这一阻尼项，Ravine 一侧的梯度可和性分析是自然的切入口。

\paragraph{噪声递减：JOTA 所在的格子。}
JOTA 的条件可以分三层：$(K_2)$ 给出 a.s.\ $O(1/(s_kt_k^2))$；
再加上 $(K_1^+)$ 与 $\sum_k t_{k+1}/t_k^2=\infty$，得到函数值的 a.s.\ $o$ 速率；
迭代点收敛还需要 $\alpha_k\in[0,1]$、$\inf_k s_k>0$ 与 $\ell^1$ 型条件 $(K_2^+)$。
其中 $\ell^1$ 条件来自对距离增量取正部的证明手法，是否必要尚无结论。

\section{JOTA 的位置与剩余问题}

JOTA 位于
\[
  \text{一般 }\alpha_k\;\times\;\text{噪声递减}\;\times\;\text{SNAG/SRAG}
\]
这一格。它的增量有三点：动量参数从“常数”或“由步长决定”推广到一般 $\alpha_k$；
噪声递减时允许不降步长，从而保住加速速率；
借助 Nesterov 与 Ravine 的对应，把结果转移到 SRAG。
持续噪声下，JOTA 与 \citet{sebbouh2021} 处于同一参数体制，速率主导阶相同，区别只在算法与参数范围。

\begin{table}[htbp]
\centering
\small
\caption{按 $\alpha$ 与噪声体制划分的已知结果与空白（$\alpha_k=1-\alpha/k$）}
\label{tab:gaps}
\begin{tabularx}{\textwidth}{@{}p{2.6cm}YYp{1.6cm}@{}}
\toprule
方向 & 已有结果 & 仍缺 & JOTA \\
\midrule
$\alpha=3$，噪声递减
& a.s.\ $O(1/k^2)$ \citep{attouch2026ravine}；
  FISTA 迭代点 a.s.\ 弱收敛 \citep{salzo2025iterates}（条件\tbv）
& a.s.\ $o(1/k^2)$；Salzo 的结果是否覆盖同一 SNAG 模型
& 部分覆盖 \\
\addlinespace
$\alpha>3$，噪声递减
& a.s.\ $o(1/k^2)$、梯度可和、迭代点收敛 \citep{attouch2026ravine}
& $(K_2^+)$ 是否必要（备用）
& 已覆盖 \\
\addlinespace
$\alpha\ge3$，持续噪声
& a.s.\ 速率任意接近 $k^{-1/2}$ \citep{attouch2026ravine}；
  SHB 的迭代点 a.s.\ 收敛 \citep{sebbouh2021,liu2022}
& SNAG/SRAG 的迭代点收敛
& 仅速率 \\
\addlinespace
$\alpha\ge3$，相对噪声
& 期望 $O(1/k^2)$ \citep{vaswani2019fast,gupta2024nesterov}；
  修正方法 \citep{yu2026shang}（a.s.\ 结论\tbv）
& a.s.\ $O/o$ 速率、迭代点收敛、对 $\alpha$ 与临界值的依赖
& 未覆盖 \\
\bottomrule
\end{tabularx}
\end{table}

据此，汇报保留三个剩余问题。

\paragraph{问题 1：相对噪声下的 a.s.\ 理论。}
设
\[
  \alpha_k=1-\frac{\alpha}{k},\qquad
  \mathbb E\big[\|e_k\|^2\mid\mathcal F_k\big]\le\nu^2\|\nabla f(y_k)\|^2 .
\]
第一式是标准 Nesterov 动量，第二式说明噪声随梯度一同消失。
问题是：在 $\nu<\nu_0(\alpha)$ 与常数步长下，是否有
\[
  f(x_k)-f^\star=o(1/k^2)\quad\text{a.s.}
\]
以及迭代点收敛？阈值 $\nu_0$ 如何依赖 $\alpha$，临界 $\alpha=3$ 时又如何？
负结果决定了 $\nu$ 必须受限，因此这里的 $\nu_0(\alpha)$ 本身就是要刻画的对象。

\paragraph{问题 2：持续噪声下 SNAG/SRAG 的迭代点收敛。}
速率已被下界锁定，但迭代点收敛是独立的问题。
SHB 已有结果 \citep{sebbouh2021,liu2022}，可作为比较对象；
而 JOTA 的迭代点结论要求 $\inf_k s_k>0$，与持续噪声下必须递减的步长相冲突。

\paragraph{问题 3：噪声递减、临界 $\alpha=3$ 下的 a.s.\ $o$ 速率。}
$(K_1^+)$ 提供的余量在 $\alpha=3$ 时消失，JOTA 的 $o$ 速率论证因此失效。
需要核对 \citet{salzo2025iterates} 是否覆盖同一随机模型，并判断其技术能否给出速率。

$\alpha<3$ 的随机情形与“噪声条件是否紧”放入备用材料，不作为本次汇报的主问题。

\bibliographystyle{plainnat}
\bibliography{refs}

\end{document}
